\section{Introduction}\label{sec:introduction}
%First paragraph
In response to the recent surge in inflation, most major central banks have started tightening their policy stance. \rev{However, because macroeconomic fundamentals, financial dynamics and cyclical positions differ across countries, monetary policies also diverge internationally: some central banks --- such as the Bank of England or the Federal Reserve (Fed) --- started tightening relatively early, while others --- most notably the European Central Bank (ECB) --- were slower to raise rates. Such divergence opens up potential interaction effects between central banks.}

These interaction effects arise because monetary policy often co-moves internationally. However, most models assume that a central bank's reaction function is driven purely by domestic factors \citep[see, e.g.,][for a recent state-of-the-art dynamic stochastic general equilibrium model that is used for policy analysis]{del2013frbny}, a rather strong assumption given the increased financial and real linkages in the world economy. Even if domestic quantities are somewhat affected by foreign activity through the trade balance, foreign monetary policy is a neglected dimension in these analyses. Therefore, our goal in this paper is to model the ECB's monetary policy as being explicitly dependent on the Fed's policy stance, given the prominent role of the latter in the global financial cycle \citep{miranda2020us}. \rev{Concretely, we estimate a Bayesian smooth transition vector autoregression (ST-VAR) for weekly euro-area financial data whose dynamic coefficients depend, through a logistic transition function, on the Fed's policy stance. We measure this stance as the deviation of the effective federal funds rate from the US natural rate of interest \citep{laubach2003measuring} --- the rate consistent with output at potential and stable inflation --- so that positive values signal a restrictive Fed. The transmission of ECB policy surprises to euro-area yields and market-based inflation expectations is thereby allowed to vary smoothly with the Fed's stance.}  


%Global modeling of spillovers
One way to capture international co-movement in monetary policy (and financial markets more generally) is through the use of large-scale, multi-country models. The literature on empirical global macro models \citep[see, among many others,][]{dees2007exploring, feldkircher2016international,georgiadis2017bi, koop2016model, burriel2018uncovering, dees2021global} typically uses large-scale models that represent the structure of the economy and are able to capture cross-country relationships. These models often assume that countries are dynamically linked and that shocks originating in one country have immediate and lagged effects not only on itself but also on other countries through spillovers. \rev{These frameworks are well suited for counterfactual and conditional forecasting exercises, in which the researcher traces out how outcomes change under alternative policy paths or modified reaction functions. Our question is different: we ask whether the \textit{propagation} of a given domestic policy shock --- the dynamic coefficients that map an ECB surprise into yields, inflation compensation, and other asset prices --- itself varies systematically with the stance of the foreign central bank. In linear multi-country models this transmission mechanism is, by construction, invariant to the foreign policy environment: a conditional forecast changes the path of the shocks, not the propagation of a given shock. Answering our question therefore requires a model in which the coefficients are an explicit function of the foreign policy stance.} Hence, the main focus of this paper is to empirically analyze how domestic monetary policy effects change in the euro area (EA) \textit{conditional} on a potentially divergent foreign policy stance of the globally most important central bank --- the Fed.

% MP spillovers incl. channels
More generally, several recent papers have focused on how monetary policy effects are attenuated by various influences. One strand of the literature addresses this question by looking at economic and policy spillovers to other economies. These empirical papers focus mainly on US monetary policy spillovers and spillbacks \citep{kim2001international, canova2005transmission, bluwstein2018beggar, crespo2019spillovers, ca2020monetary, breitenlechner2021spillbacks, ca2021making, gj2023global} and find strong effects on other economies. 

In general, the empirically identified channels through which foreign monetary policy may affect domestic macroeconomic quantities consider both real and financial markets. In this paper, we focus on the channel operating through financial markets. Countries are exposed to changing financial conditions, exchange rate dynamics, and collateral valuation effects with potential implications for bank balance sheets. Here, the US as a supplier of safe assets (Treasuries), its impact on the global financial cycle \citep{miranda2020us}, and the special role of the US dollar \citep{gopinath2021dominant} are central parts of our analysis. Thus, changes in the demand for US assets triggered by US monetary policy actions can change the demand for EA assets. As a consequence, policy actions of other central banks, especially if not aligned with the Fed's, may affect financial variables with potentially offsetting effects.

\rev{Why should the transmission of an ECB surprise depend on the Fed's stance? We emphasize three complementary mechanisms, each with a testable implication that we take to the data. First, an \textit{interest parity and carry trade channel}: when the Fed is accommodative, an unexpected ECB tightening widens transatlantic yield differentials from a compressed level, triggering portfolio flows into EA bonds. This predicts stronger short-run yield reactions --- particularly at short and medium maturities relevant for currency-hedged carry positions --- followed by a partial reversal as inflows bid up bond prices. Second, a \textit{global financial cycle and risk-taking channel} \citep{miranda2020us}: a loose Fed compresses risk premia globally and encourages leverage; in such an environment, a domestic policy surprise moves term premia and risky asset prices by more than it would in a high-premium, deleveraged environment. This predicts larger long-end and stock market responses under a loose Fed. Third, a \textit{safe asset and dominant currency channel} \citep{gopinath2021dominant}: AAA-rated EA bonds are the closest substitutes for Treasuries, so the pricing impact of shifts in relative supply and demand for safe assets depends on the Fed-driven return on the dominant safe asset. Importantly, all three mechanisms operate through the \textit{financial} side and are conceptually distinct from state dependence on the domestic business cycle, which the literature has studied extensively with mixed results \citep[see, e.g.,][]{tenreyro2016pushing}: the conditioning variable is the foreign policy stance, not domestic slack. We show below that the two are empirically distinct as well --- regime allocations implied by the Fed's stance are essentially uncorrelated with those implied by euro area cyclical conditions.} As recently shown by \cite{ca2021making}, while US monetary policy shocks have a noticeable impact on several EA macroeconomic and financial indicators, the spillover effects of EA monetary policy on the US economy are rather limited, if not absent. Interestingly, the study also found that the effects on real economic activity and, in particular, financial markets were much larger. To summarize, most empirical studies have focused on the direct effects of monetary policy surprises, i.e., whether identified exogenous shocks have a direct impact on the domestic economy. However, there is a notable lack of research on whether the effects of domestic monetary policy shocks depend on the policy stance of a foreign bank. 

\rev{We fill this gap with the ST-VAR sketched above, in which the dynamic coefficients vary with the Fed's stance as the signal variable that determines the regime.} ECB monetary policy shocks are identified through high-frequency reactions of financial markets in a narrow window around scheduled events associated with monetary policy announcements \cite{altavilla2019measuring}. \rev{This lets us track how the effects of ECB policy surprises change with the Fed's stance across the term structure.}

Our analysis looks at the financial side of the economy and includes weekly data on EA bond yields of different maturities and inflation-linked swaps to proxy marked-based inflation expectations. Since the number of different maturities is large, we rely on a term structure model in the Nelson-Siegel tradition \citep{DIEBOLD2006337}.  This allows us to capture the dynamic responses of real interest rates over a wide range of maturities. Compared to existing studies, we provide a comprehensive analysis that includes not only the full term structure but also explicit conditions on the Fed's stance in a high-frequency setting. This helps to shed light on how international linkages shape financial market reactions to monetary policy. 



%Our model is complemented by including  inflation-linked swaps as a proxy for market-based inflation expectations. In addition, we use traditional term structure models in the Nelson-Siegel tradition \citep{DIEBOLD2006337} to obtain the full term structure of interest rates and market-based inflation expectations. 



%Thus, we focus on the transmission of the ECB's policy through financial markets at a \textit{weekly} frequency. In addition, we choose a nonlinear model to account for the monetary policy stance of the US Federal Reserve. We use a smooth transition vector autoregression (ST-VAR) model that takes the Fed's policy stance as the signal variable determining the regime. This allows us to flexibly track how the effects of ECB policy surprises change based on the Fed's stance. Our analysis looks at the financial side of the economy and uses inflation-linked swaps to augment our approach with a proxy for market-based inflation expectations. In addition, we use traditional term structure models in the Nelson-Siegel tradition \citep{DIEBOLD2006337} to obtain the full term structure of interest rates and market-based inflation expectations. This allows us to capture the dynamic responses of real interest rates over a wide range of maturities. Compared to existing studies, we provide a comprehensive analysis that includes not only the full term structure but also explicit conditions on the Fed's stance in a high-frequency setting. This helps to shed light on how international linkages shape financial market reactions to monetary policy. 

\sug{[Suggested calibration for the revision. A model-free look at the data --- the pass-through of ECB surprises across the term structure (\autoref{sec:passthrough}) --- confirms that our instrument transmits cleanly to short-maturity nominal and real yields, but is too imprecise at weekly frequency to establish whether that transmission is \emph{state-dependent}. The state-dependence claims below should therefore remain explicitly attributed to the ST-VAR (as they already are), so that no headline claim rests on the simple reduced-form regressions. We suggest adding one sentence to this effect here, mirroring the abstract.]}
The empirical results indicate that the Fed had a rather restrictive stance until the onset of the financial crisis and after 2016. The period in between, however, was characterized by expansionary periods punctuated by several restrictive and transitory policy regimes. With respect to the contractionary monetary surprises in the EA, our findings suggest that the dynamic responses of EA financial markets strongly depend on the respective Fed stance, providing appreciable evidence of nonlinear interaction effects. Moreover, the relationship between impulse responses and US monetary policy is heterogeneous with respect to shape and magnitude over the term structure. At the short end of the curve, government bond yields increase, with a consecutive decrease only in the case of an expansionary Fed stance. The longer end of the curve shows declining yields after one week when the Fed is hawkish. Inflation-linked swaps initially fall, but then rise after one week when the Fed is expansionary. However, the differences in responses between  expansionary and contractionary US policy phases are not significant over the entire term structure.

\rev{We subject these findings to an extensive battery of checks that speak directly to alternative interpretations. First, to separate the Fed's stance from common global forces --- such as a synchronized global inflation cycle in which many central banks tighten simultaneously --- we augment the model with global variables (oil prices, the broad US dollar index, and the VIX) and, alternatively, purge the stance measure of global factors before using it as the transition variable. The results are virtually unchanged. Second, we document that the regime allocation implied by the Fed's stance is essentially uncorrelated with the one implied by a euro area business cycle indicator, and the state dependence survives when tightening and easing surprises are considered separately. This alleviates the concern that the Fed's stance merely proxies the domestic cycle or the effective lower bound period. Third, our sample runs through the end of 2025 and thus covers the post-pandemic global tightening --- the first pronounced Fed tightening outside a recessionary environment in the sample. Far from overturning the result, this episode sharpens it: euro area policy surprises move financial markets more strongly when the Fed is accommodative, and the recent restrictive-Fed observations are what pin down the response at the tight end of the stance. Finally, we show that the fast mean reversion of the estimated responses reflects the fact that the model is specified in first differences: the implied \textit{level} effects on yields are persistent.}

One key advantage of our approach is that it allows us to examine how the term structure of \textit{real} rates changes in response to EA monetary policy shocks conditional on the prevailing US monetary policy regime. This analysis reveals that  restrictive ECB surprises trigger stronger positive reactions if the Fed is hawkish as opposed to a situation where monetary policy diverges.


The rest of the paper is organized as follows. In the next section, we introduce our econometric model, while in \autoref{sec:Data} we discuss the dataset, perform univariate analyses to set the stage, and show the model-implied US monetary policy regimes.  \autoref{sec:results} presents the main results based on our ST-VAR as well as a robustness exercise. The last section summarizes and concludes the paper.


%The first channel, the aggregate demand channel, operates through international trade and becomes more important with increasing openness. In a nutshell, a domestic contractionary monetary policy shock leads to a deterioration of domestic demand through lower consumption and investments, such that also imports (and thus exports of a foreign country) decline.%\textbf{CITE} %Expand and provide empirical support. 
%A second channel concerns the relative price between foreign and domestic goods and the influence of the exchange rate. However, this channel depends on various factors, like  price stickiness, the currencies involved or the degree of the exchange rate pass-through and is thus very specific to the countries under scrutiny. For our application, the dominant role of the US-dollar needs to be considered, as it has major implications for the international transmission of shocks as pointed out by \cite{gopinath2021dominant}.
%Finally, the last channel operates via financial markets and foreign countries are exposed to changing financial conditions, exchange rate dynamics and valuation effects of collaterals with potential effects on balance sheets. For our present application, we take a closer look on this last channel. Especially, the role of the US as a supplier of safe assets (treasuries) and its sensitivity to monetary policy will be a central part of our analysis. Changes in the demand for US assets triggered by tighter US monetary policy may thus increase the demand for EA assets. Hence, policy actions from both parties may affect the EA yield curve and other financial variables with potentially offsetting effects.
